\chapter{Desarrollo}
\label{cap:desarrollo}

En este capítulo se presenta la implementación práctica del proceso, comenzando por la reproducción del circuito descrito en el paper de referencia~\cite{arian2025qhrp} y continuando con la implementación del circuito que realice el proceso por etapas, de tal manera, que se facilite su trazabilidad.

\section{Circuito cuántico extraído del paper de referencia}
\label{sec:circuito_extraido_paper}

Como parte de la validación metodológica, se construye una representación esquemática del circuito de referencia descrito en el artículo~\cite{arian2025qhrp} para el caso de seis qubits. La figura~\ref{FIG:CIRCUITO_PAPER_QHRP} ilustra la estructura general del mapa de características tal y como se presenta en dicho artículo. La secuencia incluye: (i) compuertas Hadamard en todos los qubits, para cambiar a la distribución equiprobable de estados; (ii) tres capas de codificación con rotaciones $R_y$, $R_x$ y de nuevo $R_y$, que parametrizan en los ángulos el mapa de características; (iii) capas de entrelazamiento intercaladas generadas con puertas CNOT; y (iv) una etapa final de medición en la base computacional. La última etapa de medición se incluye para su ejecución en ordenadores reales, siendo innecesaria cuando se utiliza el simulador de vector de estados.

Por otro lado, indicar que la puerta CNOT, es una opción para generar entrelazamiento, pero las posibilidades de equivalencia de puertas permiten otras alternativas de generarlo con puertas como \texttt{CZ}, \texttt{SWAP} e \texttt{iSWAP}, y con ello obtener relaciones no lineales entre las características que se representan en el circuito.

La representación de la figura~\ref{FIG:CIRCUITO_PAPER_QHRP} es la base del diseño del ansatz y el punto de partida para facilitar la comparación entre la formulación teórica y su implementación práctica.

Para evitar confusiones de interpretación, entre el circuito del paper de referencia y el utilizado en el código experimental se detallan dos circuitos complementarios:
\begin{itemize}
\item \textbf{Circuito de referencia (ver Figura~\ref{FIG:CIRCUITO_PAPER_QHRP}):} replica la narrativa del paper de referencia~\cite{arian2025qhrp} y usa CNOT como caso ilustrativo de entrelazamiento.
\item \textbf{Circuito del proceso experimental \gls{QHRP}:} usa la secuencia \texttt{H} -- \texttt{T} -- \texttt{RY/RX} con capas de \texttt{iSWAP} en pares predefinidos, consistente con el módulo de simulación que genera las matrices de densidad.
\end{itemize}

Esta decisión metodológica permite, simultáneamente, (i) mantener fidelidad documental con la presentación del paper de referencia~\cite{arian2025qhrp} y (ii) preservar coherencia con la implementación reproducible que produce los resultados reportados. Ambos circuitos son esencialmente equivalentes desde el punto de vista del mapeo de características cuánticas, pero la sustitución de CNOT por iSWAP modifica la geometría del circuito implementado. Puede haber cambios adicionales cuando se produce la traspilación para su ejecución en un ordenador cuántico real.

\begin{figure}[Circuito cuántico reproducido del paper de referencia]{FIG:CIRCUITO_PAPER_QHRP}
{Representación esquemática del circuito cuántico de 6 qubits basada en el esquema del artículo de referencia~\cite{arian2025qhrp}.}
\image{15cm}{}{fig1-circuito-paper-qhrp}
\end{figure}

El circuito confirma que cada una de las 6 características del activo se codifica en un qubit y se correlaciona mediante capas de entrelazamiento. El mapa de características define un espacio de Hilbert de dimensión $2^6=64$ para 6 qubits, donde la información queda codificada en la geometría de puertas y los parámetros de las mismas que definen las correlaciones cuánticas.

\subsection{Codificación en Qiskit de los circuitos}
\label{subsec:codificacion_qiskit_circuitos}

La codificación y documentación en el repositorio git se organiza en 2 niveles:
\begin{itemize}
\item \textbf{Circuito de referencia (Figura~\ref{FIG:CIRCUITO_PAPER_QHRP}):} notebook que construye el esquema del circuito de referencia de 6 qubits y permite visualizar su estructura con las capas de codificación y entrelazamiento.
\item \textbf{Circuito del proceso experimental \gls{QHRP}:} módulo principal de QHRP (archivo \texttt{src/qhrp\_figure3\_4/src/quantum\_hrp.py}), donde se define el ansatz usado para calcular estados, matrices de densidad promedio y distancias de Frobenius.
\end{itemize}

La trazabilidad queda garantizada porque el notebook documenta la arquitectura de referencia y el módulo del proceso experimental e implementa la configuración que da lugar a todos los resultados reportados. De este modo, es posible reproducir tanto la figura de referencia como los resultados cuantitativos del trabajo bajo el proceso experimental utilizado.

\subsection{Reproducibilidad del entorno, versiones y trazabilidad}
\label{subsec:reproducibilidad_entorno}

Con el objetivo de facilitar la replicabilidad de los resultados, el proyecto mantiene dos ficheros de dependencias:
\begin{itemize}
\item \texttt{requirements\_paper\_qiskit.txt}: entorno compatible con la versión usada para reproducir el proceso del paper de referencia~\cite{arian2025qhrp}.
\item \texttt{requirements\_act\_qiskit.txt}: entorno actualizado para pruebas recientes y extensiones.
\end{itemize}

En particular, se documentan explícitamente las librerías clave de cálculo científico, de visualización y de computación cuántica utilizadas en los experimentos.

Para cálculo científico se emplean \texttt{numpy}, \texttt{scipy} y \texttt{pandas}; para visualización, \texttt{matplotlib} y \texttt{seaborn}; y para computación cuántica, \texttt{qiskit} y módulos asociados.

Para ejecutar el código en backend real de IBM, es necesario instalar \texttt{qiskit-ibm-runtime} y proporcionar credenciales válidas de IBM Quantum.

Para facilitar los experimentos, se han utilizado 3 notebooks principales, para diferenciar las pruebas en simulación, las pruebas con hardware real y los experimentos con 5 qubits:
\begin{itemize}
\item Notebook principal con resultados obtenidos de simulación.
\item Notebook independiente para prueba en hardware real siguiendo el mismo flujo de proceso que la simulación.
\item Notebook específico para experimentos con 5 qubits.
\end{itemize}

Esta organización permite distinguir con claridad resultados en simulación ideal, resultados en ejecución hardware y resultados de comparación con el paper de referencia~\cite{arian2025qhrp}, evitando ambigüedades en la interpretación.

\section{Cálculo de distancias, estructura jerárquica y asignación de pesos}
\label{sec:calculo_distancias_estructura}

Una vez obtenidas las matrices densidad promedio, se calcula la distancia de Frobenius entre todos los pares de activos. Estas distancias se almacenan en una matriz simétrica $D$.

El orden de los activos se determina mediante técnicas de clustering jerárquico, siguiendo la metodología \gls{HRP} original. Este procedimiento permite identificar la estructura de grupos de activos basándose en la proximidad de sus estados cuánticos. Finalmente, el sistema ejecuta el algoritmo de bisección recursiva clásico sobre la matriz de covarianza para calcular los pesos definitivos de la cartera, garantizando así la paridad de riesgo basada en la nueva topología cuántica.

En la comparación experimental, \gls{HRP} clásico y \gls{KHRP} emplean enlace simple, mientras que la ordenación cuántica evaluada utiliza Ward sobre las distancias de Frobenius. La fase de asignación de pesos sigue siendo la misma bisección recursiva clásica, de modo que la diferencia metodológica queda concentrada en el criterio de ordenamiento.

\section{Generación de visualizaciones}
\label{sec:generacion_visualizaciones}

Las visualizaciones se generan mediante mapas de calor dispobibles en la librería \texttt{seaborn}, que representan las matrices de distancia cuántica. Se compara la matriz original desordenada, con la matriz reordenada mediante la jerarquía cuántica, lo que permite observar visualmente la aparición de bloques cuasi-diagonales que indican agrupaciones de activos similares.

\section{Descarga y preparación de datos financieros}
\label{sec:descarga_preparacion_datos}

La cartera de activos, el procedimiento de descarga de precios desde Yahoo Finance y las etapas de preprocesado y construcción de características financieras se describen en detalle en el Apéndice~\ref{ap:datos}. Allí se especifican los 170 activos del universo original, el período 2005--2025, los pasos de limpieza y cálculo de retornos, y las seis características calculadas mediante ventanas móviles que constituyen la entrada al mapa de características cuántico. El tensor resultante de dimensiones $(N, T, 6)$ se almacena en formato binario (\texttt{.npy}) para desacoplar la preparación de datos del análisis posterior y facilitar la reproducibilidad.
